\chapter{Sheaves of Modules}

\section{\texorpdfstring{\(\OO_X\)}{OX}-Modules}

\begin{definition}
Let \(X\) be a scheme.  An \emph{\(\OO_X\)-module} is a sheaf \(\FF\) of
abelian groups on \(X\) such that each \(\FF(U)\) is an \(\OO_X(U)\)-module and
restriction maps are compatible with scalar multiplication.  A morphism of
\(\OO_X\)-modules is a morphism of sheaves of abelian groups that is
\(\OO_X\)-linear on every open set.
\end{definition}

\begin{definition}
If \(\FF\) and \(\GG\) are \(\OO_X\)-modules, their tensor product is the
sheafification of the presheaf
\[
    U\longmapsto \FF(U)\tensor_{\OO_X(U)}\GG(U).
\]
It is denoted \(\FF\tensor_{\OO_X}\GG\).  The sheaf
\(\shHom_{\OO_X}(\FF,\GG)\) is defined by
\[
    U\longmapsto \Hom_{\OO_U}(\FF|_U,\GG|_U).
\]
\end{definition}

\begin{proposition}
For every point \(x\in X\),
\[
    (\FF\tensor_{\OO_X}\GG)_x
    \simeq
    \FF_x\tensor_{\OO_{X,x}}\GG_x.
\]
\end{proposition}

\begin{proof}
Taking the stalk of the presheaf tensor product gives the filtered colimit of
\(\FF(U)\tensor_{\OO_X(U)}\GG(U)\) over neighbourhoods \(U\) of \(x\).
Filtered colimits commute with tensor products of modules, and the filtered
colimit of the rings \(\OO_X(U)\) is \(\OO_{X,x}\).  Sheafification does not
change stalks, so the displayed isomorphism follows.
\end{proof}

\begin{definition}
Let \(f:X\to Y\) be a morphism of schemes.  If \(\GG\) is an \(\OO_Y\)-module,
its \emph{pullback} is
\[
    f^*\GG=f^{-1}\GG\tensor_{f^{-1}\OO_Y}\OO_X,
\]
where \(\OO_X\) is an \(f^{-1}\OO_Y\)-module through \(f^\#\).
\end{definition}

\begin{proposition}
For \(x\in X\), there is a natural isomorphism
\[
    (f^*\GG)_x\simeq \GG_{f(x)}\tensor_{\OO_{Y,f(x)}}\OO_{X,x}.
\]
\end{proposition}

\begin{proof}
The stalk of \(f^{-1}\GG\) at \(x\) is \(\GG_{f(x)}\).  Taking stalks in the
definition of \(f^*\GG\) and using the preceding proposition gives the stated
isomorphism.
\end{proof}

\section{Modules on Affine Schemes}

\begin{definition}
Let \(A\) be a ring and \(M\) an \(A\)-module.  On \(X=\Spec A\), define
\(\widetilde{M}\) by
\[
    \widetilde{M}(D(f))=M_f
\]
on basic opens, with the evident localization maps.  Extending by the same
local fraction construction as for \(\OO_X\) gives an \(\OO_X\)-module.
\end{definition}

\begin{proposition}
For \(X=\Spec A\), \(M\) an \(A\)-module, and \(f\in A\), the natural map
\[
    M_f\longrightarrow \widetilde{M}(D(f))
\]
is an isomorphism.  In particular \(\Gamma(X,\widetilde{M})\simeq M\), and
\((\widetilde M)_{\mathfrak p}\simeq M_{\mathfrak p}\).
\end{proposition}

\begin{proof}
The proof is identical to the proof for the structure sheaf, with elements of
\(M\) in place of elements of \(A\).  Injectivity is checked after localizing at
all primes of \(A_f\).  For surjectivity, a section over \(D(f)\) is locally of
the form \(m_i/g_i^{n_i}\) on a finite principal cover \(D(g_i)\).  After
clearing denominators and using a relation \(f^N=\sum_i b_i g_i\), the element
\(\sum_i b_i m_i/f^N\in M_f\) restricts to the given local sections.  The stalk
statement follows by taking the filtered colimit over \(D(f)\ni\mathfrak p\).
\end{proof}

\begin{proposition}
The functor
\[
    M\longmapsto \widetilde M
\]
from \(A\)-modules to \(\OO_{\Spec A}\)-modules is exact and commutes with
arbitrary direct sums, tensor products, and localization.
\end{proposition}

\begin{proof}
Exactness can be checked on stalks.  For an exact sequence of \(A\)-modules
\[
    M'\longrightarrow M\longrightarrow M'',
\]
the induced sequence on the stalk at \(\mathfrak p\) is
\[
    M'_{\mathfrak p}\longrightarrow M_{\mathfrak p}\longrightarrow
    M''_{\mathfrak p},
\]
which is exact because localization is exact.  Hence the sequence of associated
sheaves is exact.  Direct sums and tensor products are also checked on basic
opens:
\[
    \left(\bigoplus_i M_i\right)_f\simeq\bigoplus_i (M_i)_f,\qquad
    (M\tensor_A N)_f\simeq M_f\tensor_{A_f}N_f.
\]
Localization is the identity
\(\widetilde M|_{D(f)}\simeq \widetilde{M_f}\) under
\(D(f)\simeq\Spec A_f\).
\end{proof}

\begin{proposition}
Let \(A\to B\) be a ring homomorphism and let
\(f:\Spec B\to\Spec A\) be the induced morphism.  For every \(A\)-module \(M\),
\[
    f^*\widetilde M\simeq \widetilde{M\tensor_A B}.
\]
\end{proposition}

\begin{proof}
It suffices to compare stalks at a prime \(\mathfrak q\subset B\), with
\(\mathfrak p=\mathfrak q\cap A\).  The stalk of \(f^*\widetilde M\) is
\[
    M_{\mathfrak p}\tensor_{A_{\mathfrak p}}B_{\mathfrak q},
\]
and the stalk of \(\widetilde{M\tensor_A B}\) is
\[
    (M\tensor_A B)_{\mathfrak q}.
\]
The canonical map between these two modules is an isomorphism by the universal
property of localization.  Since a morphism of sheaves is an isomorphism when
it is an isomorphism on stalks, the result follows.
\end{proof}

\section{Quasi-Coherent and Coherent Sheaves}

\begin{definition}
An \(\OO_X\)-module \(\FF\) is \emph{quasi-coherent} if every point of \(X\) has
an affine open neighbourhood \(U=\Spec A\) such that
\(\FF|_U\simeq\widetilde M\) for some \(A\)-module \(M\).  It is
\emph{coherent} if \(X\) is locally noetherian and every point has an affine
open neighbourhood \(U=\Spec A\) such that
\(\FF|_U\simeq\widetilde M\) for a finitely generated \(A\)-module \(M\).
\end{definition}

\begin{theorem}
Let \(X=\Spec A\).  The functor \(M\mapsto\widetilde M\) gives an equivalence
between \(A\)-modules and quasi-coherent \(\OO_X\)-modules.  A quasi-coherent
sheaf \(\FF\) corresponds to the module \(\Gamma(X,\FF)\).
\end{theorem}

\begin{proof}
We already know that \(\Gamma(X,\widetilde M)=M\), so the functor is faithful
and injective on objects up to isomorphism.  Let \(\FF\) be quasi-coherent.
Cover \(X\) by basic opens \(D(f_i)\) such that
\(\FF|_{D(f_i)}\simeq\widetilde{M_i}\).  We prove that the canonical map
\[
    \widetilde{\Gamma(X,\FF)}\longrightarrow \FF
\]
is an isomorphism.  It is enough to check on the basic opens \(D(f_i)\).  On
\(D(f_i)\), the global sections of \(\FF\) localize to
\(\Gamma(D(f_i),\FF)\): if \(s\in\Gamma(D(f_i),\FF)\), quasi-compactness of
\(D(f_i)\) and the sheaf condition allow \(s\) to be represented after
localization by a global section divided by a power of \(f_i\).  This is the
same clearing-denominators argument used for \(\widetilde M\).  Thus
\[
    \Gamma(X,\FF)_{f_i}\simeq\Gamma(D(f_i),\FF),
\]
and the canonical map is an isomorphism on the cover.  Hence \(\FF\) is
\(\widetilde{\Gamma(X,\FF)}\).

For morphisms, if \(\FF=\widetilde M\) and \(\GG=\widetilde N\), every
\(\OO_X\)-linear morphism \(\FF\to\GG\) is determined by the induced
\(A\)-linear map on global sections \(M\to N\), and every such map sheafifies.
Therefore the functor is full and faithful.
\end{proof}

\begin{proposition}
On a scheme \(X\), kernels, cokernels, images, extensions, direct sums, and
tensor products of quasi-coherent sheaves are quasi-coherent.  If \(X\) is
locally noetherian, the same statement holds for coherent sheaves, with finite
direct sums.
\end{proposition}

\begin{proof}
All assertions are local on affine opens.  On an affine open \(\Spec A\),
quasi-coherent sheaves correspond to \(A\)-modules.  Kernels, cokernels,
images, extensions, direct sums, and tensor products of modules are modules,
and the associated sheaf functor is exact and commutes with sums and tensor
products.  Thus the quasi-coherent case follows.

If \(A\) is noetherian, submodules, quotients, extensions, finite direct sums,
and tensor products of finitely generated \(A\)-modules are finitely generated.
Therefore the coherent case follows from the same affine reduction.
\end{proof}

\begin{proposition}
Let \(f:X\to Y\) be a morphism of schemes.  The pullback of a quasi-coherent
\(\OO_Y\)-module is quasi-coherent.  If \(X\) and \(Y\) are locally noetherian
and \(f\) is locally of finite type, the pullback of a coherent sheaf is
coherent.
\end{proposition}

\begin{proof}
The question is local on \(X\) and \(Y\).  On affine opens the assertion is the
formula
\[
    f^*\widetilde M\simeq \widetilde{M\tensor_A B}
\]
for a ring map \(A\to B\).  This proves quasi-coherence.  If \(M\) is finitely
generated and \(B\) is a finitely generated \(A\)-algebra, then
\(M\tensor_A B\) is finitely generated as a \(B\)-module.  This proves
coherence in the stated noetherian finite type situation.
\end{proof}

\section{Invertible Sheaves}

\begin{definition}
An \(\OO_X\)-module \(\LL\) is \emph{invertible} if every point of \(X\) has an
open neighbourhood \(U\) such that \(\LL|_U\simeq \OO_U\).  An invertible sheaf
is also called a \emph{line bundle}.  The tensor product makes the set of
isomorphism classes of invertible sheaves into an abelian group, the
\emph{Picard group} \(\Pic(X)\).
\end{definition}

\begin{proposition}
If \(\LL\) is invertible, then
\[
    \LL^\vee=\shHom_{\OO_X}(\LL,\OO_X)
\]
is inverse to \(\LL\) in \(\Pic(X)\); that is,
\(\LL\tensor_{\OO_X}\LL^\vee\simeq\OO_X\).
\end{proposition}

\begin{proof}
The evaluation morphism
\[
    \LL\tensor_{\OO_X}\LL^\vee\longrightarrow\OO_X
\]
is local on \(X\).  On an open set where \(\LL\simeq\OO_X\), it becomes the
standard isomorphism
\[
    \OO_X\tensor_{\OO_X}\shHom_{\OO_X}(\OO_X,\OO_X)\simeq\OO_X.
\]
Since being an isomorphism can be checked locally, the evaluation morphism is
an isomorphism globally.
\end{proof}

\begin{proposition}
Let \(\LL\) be an invertible sheaf and let
\[
    0\longrightarrow \FF'\longrightarrow \FF\longrightarrow \FF''
    \longrightarrow 0
\]
be an exact sequence of \(\OO_X\)-modules.  Then
\[
    0\longrightarrow \FF'\tensor\LL\longrightarrow \FF\tensor\LL
    \longrightarrow \FF''\tensor\LL\longrightarrow 0
\]
is exact.
\end{proposition}

\begin{proof}
Exactness is local on \(X\).  On an open set where \(\LL\simeq\OO_X\), the
sequence after tensoring with \(\LL\) is identified with the original exact
sequence.  These local exactness statements imply exactness of the sheaf
sequence because exactness can be checked on stalks.
\end{proof}

\section{Finite Presentation and Support}

\begin{definition}
An \(\OO_X\)-module \(\FF\) is \emph{of finite type} if every point has an open
neighbourhood \(U\) on which there is a surjection
\[
    \OO_U^{\oplus n}\longrightarrow \FF|_U
\]
for some \(n\).  It is \emph{of finite presentation} if locally there is an
exact sequence
\[
    \OO_U^{\oplus m}\longrightarrow \OO_U^{\oplus n}
    \longrightarrow \FF|_U\longrightarrow0.
\]
\end{definition}

\begin{proposition}
Let \(X\) be locally noetherian.  An \(\OO_X\)-module is coherent if and only
if it is quasi-coherent and of finite type.  Equivalently, it is
quasi-coherent and locally of finite presentation.
\end{proposition}

\begin{proof}
On an affine open \(U=\Spec A\) with \(A\) noetherian, a quasi-coherent sheaf is
\(\widetilde M\).  It is of finite type precisely when \(M\) is a finitely
generated \(A\)-module.  Since \(A\) is noetherian, every finitely generated
module is finitely presented.  Thus the three conditions agree locally on an
affine noetherian cover.
\end{proof}

\begin{definition}
For an \(\OO_X\)-module \(\FF\), the \emph{support} is
\[
    \Supp(\FF)=\{x\in X:\FF_x\ne0\}.
\]
\end{definition}

\begin{proposition}
If \(X\) is locally noetherian and \(\FF\) is coherent, then \(\Supp(\FF)\) is a
closed subset of \(X\).
\end{proposition}

\begin{proof}
The assertion is local on \(X\).  Let \(X=\Spec A\) and \(\FF=\widetilde M\)
with \(M\) finite.  Then
\[
    \Supp(\FF)=\{\mathfrak p:M_{\mathfrak p}\ne0\}=V(\Ann(M)).
\]
Indeed, \(M_{\mathfrak p}=0\) if and only if some element outside
\(\mathfrak p\) annihilates all of a finite set of generators of \(M\), which is
equivalent to \(\Ann(M)\not\subset\mathfrak p\).  Thus the support is closed.
\end{proof}

\begin{proposition}
Let
\[
    0\to\FF'\to\FF\to\FF''\to0
\]
be an exact sequence of coherent sheaves on a locally noetherian scheme.  Then
\[
    \Supp(\FF)=\Supp(\FF')\cup\Supp(\FF'').
\]
\end{proposition}

\begin{proof}
At a point \(x\), the stalk sequence
\[
    0\to\FF'_x\to\FF_x\to\FF''_x\to0
\]
is exact.  The middle term is zero if and only if both the submodule and the
quotient are zero.  Hence \(x\) lies in the support of \(\FF\) exactly when it
lies in the support of \(\FF'\) or \(\FF''\).
\end{proof}

\section{Locally Free Sheaves and Determinants}

\begin{definition}
An \(\OO_X\)-module \(\mathcal E\) is \emph{locally free of rank \(r\)} if
every point has an open neighbourhood \(U\) such that
\[
    \mathcal E|_U\simeq \OO_U^{\oplus r}.
\]
A locally free sheaf of finite rank is also called a vector bundle.
\end{definition}

\begin{proposition}
Let \(\mathcal E\) be locally free of rank \(r\).  Then
\[
    \mathcal E^\vee=\shHom_{\OO_X}(\mathcal E,\OO_X)
\]
is locally free of rank \(r\), and the natural morphism
\[
    \mathcal E\longrightarrow \mathcal E^{\vee\vee}
\]
is an isomorphism.
\end{proposition}

\begin{proof}
The statement is local on \(X\).  On an open set where
\(\mathcal E\simeq\OO_X^{\oplus r}\), its dual is also
\(\OO_X^{\oplus r}\), and the double-dual map is the usual identity map on a
free module of rank \(r\).  These local isomorphisms glue.
\end{proof}

\begin{definition}
For a locally free sheaf \(\mathcal E\) of rank \(r\), its
\emph{determinant} is
\[
    \det(\mathcal E)=\bigwedge^r\mathcal E.
\]
It is an invertible sheaf.
\end{definition}

\begin{proposition}
Given an exact sequence of locally free sheaves of finite rank
\[
    0\to\mathcal E'\to\mathcal E\to\mathcal E''\to0,
\]
there is a canonical isomorphism
\[
    \det(\mathcal E)\simeq
    \det(\mathcal E')\tensor\det(\mathcal E'').
\]
\end{proposition}

\begin{proof}
The assertion is local.  Over a local ring, a surjection from a finite free
module to a finite free module splits after choosing lifts of a basis of the
quotient.  Thus locally
\(\mathcal E\simeq\mathcal E'\oplus\mathcal E''\).  For free modules, the top
exterior power of a direct sum has the canonical decomposition
\[
    \bigwedge^{r'+r''}(E'\oplus E'')
    \simeq
    \bigwedge^{r'}E'\tensor \bigwedge^{r''}E''.
\]
The construction is independent of the local splitting because it is
characterized by the alternating product of ordered bases, so the local
isomorphisms glue.
\end{proof}

\section{Quasi-Coherent Sheaves on Proj}

\begin{definition}
Let \(S=\bigoplus_{n\ge0}S_n\) be a graded ring and let \(M\) be a graded
\(S\)-module.  The sheaf \(\widetilde M\) on \(X=\Proj S\) is defined by
\[
    \widetilde M(D_+(f))=M_{(f)},
\]
where \(M_{(f)}\) is the degree-zero part of the graded localization \(M_f\).
\end{definition}

\begin{proposition}
For every homogeneous \(f\in S_+\), the restriction of \(\widetilde M\) to
\(D_+(f)\simeq\Spec S_{(f)}\) is the affine sheaf associated to the
\(S_{(f)}\)-module \(M_{(f)}\).  In particular \(\widetilde M\) is
quasi-coherent, and it is coherent if \(S\) is noetherian and \(M\) is finitely
generated.
\end{proposition}

\begin{proof}
On a smaller basic open \(D_+(fg)\subset D_+(f)\), sections are obtained by
localizing \(M_{(f)}\) at the degree-zero element represented by a suitable
homogeneous power of \(g/f\).  This is exactly the localization rule for the
affine sheaf associated to \(M_{(f)}\).  Hence the restriction is affine
quasi-coherent.  Coherence follows because localization of a finitely generated
module over a noetherian ring is finitely generated.
\end{proof}

\begin{definition}
For an integer \(n\), let \(M(n)\) be the graded module with
\[
    M(n)_d=M_{n+d}.
\]
On \(X=\Proj S\), set
\[
    \OO_X(n)=\widetilde{S(n)}.
\]
The sheaf \(\OO_X(1)\) is the \emph{Serre twisting sheaf}.
\end{definition}

\begin{proposition}
If \(S=A[x_0,\ldots,x_n]\) with the standard grading, then
\(\OO_{\PP^n_A}(1)\) is invertible.
\end{proposition}

\begin{proof}
On \(D_+(x_i)\), multiplication by \(x_i\) identifies \(S(1)_{(x_i)}\) with
\(S_{(x_i)}\).  Hence \(\OO(1)\) is free of rank one on each standard open.
On overlaps the transition function is \(x_i/x_j\), a unit there.  Thus the
local trivializations glue to an invertible sheaf.
\end{proof}

\begin{proposition}
For every quasi-coherent sheaf \(\FF\) on \(X=\Proj S\), define
\[
    \FF(n)=\FF\tensor_{\OO_X}\OO_X(n).
\]
If \(\FF=\widetilde M\), then \(\FF(n)\simeq\widetilde{M(n)}\).
\end{proposition}

\begin{proof}
It suffices to check on the basic opens \(D_+(f)\).  There,
\[
    \FF(n)(D_+(f))
    \simeq
    M_{(f)}\tensor_{S_{(f)}} S(n)_{(f)}
    \simeq
    M(n)_{(f)}.
\]
These isomorphisms commute with restriction maps, so they glue.
\end{proof}

\section{Ideal Sheaves and Closed Subschemes}

\begin{definition}
Let \(i:Z\hookrightarrow X\) be a closed subscheme.  Its \emph{ideal sheaf} is
\[
    \II_Z=\ker(\OO_X\to i_*\OO_Z).
\]
\end{definition}

\begin{proposition}
The construction \(Z\mapsto\II_Z\) gives a bijection between closed subschemes
of \(X\) and quasi-coherent ideal sheaves \(\II\subset\OO_X\).
\end{proposition}

\begin{proof}
The assertion is local on \(X\).  On \(X=\Spec A\), closed subschemes are
\(\Spec(A/I)\), and their ideal sheaves are \(\widetilde I\).  Conversely, a
quasi-coherent ideal sheaf is \(\widetilde I\) on an affine open and gives
\(\Spec(A/I)\).  The affine constructions are compatible on overlaps by
localization, hence glue.
\end{proof}

\begin{proposition}
Let \(i:Z\hookrightarrow X\) be a closed immersion.  The direct image functor
\[
    i_*:\QCoh(Z)\longrightarrow \QCoh(X)
\]
is exact and fully faithful, and its essential image consists of
quasi-coherent sheaves \(\FF\) on \(X\) annihilated locally by the ideal sheaf
\(\II_Z\).
\end{proposition}

\begin{proof}
The statement is affine local.  Let \(X=\Spec A\), \(Z=\Spec(A/I)\).  A
quasi-coherent sheaf on \(Z\) is an \(A/I\)-module \(M\).  Its direct image is
the same abelian group viewed as an \(A\)-module with \(I\) acting by zero.
This restriction-of-scalars functor from \(A/I\)-modules to \(A\)-modules is
exact and fully faithful onto the subcategory of \(A\)-modules killed by
\(I\).  Sheafifying gives the result.
\end{proof}

\section{Exercises Used Later}

\begin{exercise}[Tensor products and direct sums]\label{ex:tensor-direct-sums}
Let \(A\) be a ring, \(\{M_i\}\) and \(\{N_j\}\) families of \(A\)-modules.
Show that
\[
    \left(\bigoplus_i M_i\right)\tensor_A
    \left(\bigoplus_j N_j\right)
    \simeq
    \bigoplus_{i,j}(M_i\tensor_A N_j).
\]
\end{exercise}

\begin{solution}
For any \(A\)-module \(P\), bilinear maps
\[
    \left(\bigoplus_i M_i\right)\times
    \left(\bigoplus_j N_j\right)\to P
\]
are the same as families of bilinear maps \(M_i\times N_j\to P\), because
each element of a direct sum has finite support.  By the universal property of
tensor product, these are the same as homomorphisms
\(\bigoplus_{i,j}(M_i\tensor_A N_j)\to P\).  The representing objects are
therefore canonically isomorphic.
\end{solution}

\begin{exercise}[Twisting transitions]\label{ex:twisting-transitions}
Compute the transition functions of \(\OO_{\PP^n_A}(m)\) on the standard cover
\(D_+(x_i)\).
\end{exercise}

\begin{solution}
The sheaf \(\OO(1)\) is trivialized on \(D_+(x_i)\) by multiplication with
\(x_i\).  On \(D_+(x_i)\cap D_+(x_j)\), the change from the \(j\)-trivialization
to the \(i\)-trivialization is multiplication by \(x_j/x_i\).  Therefore the
transition function for \(\OO(m)\) is \((x_j/x_i)^m\), with the evident
interpretation for negative \(m\).
\end{solution}

\begin{exercise}[Support of a quotient]\label{ex:support-quotient}
Let \(X\) be locally noetherian and let \(\FF\to\GG\) be a surjection of
coherent sheaves.  Show that \(\Supp(\GG)\subset\Supp(\FF)\).
\end{exercise}

\begin{solution}
At every point \(x\), the map \(\FF_x\to\GG_x\) is surjective.  If
\(\FF_x=0\), then \(\GG_x=0\).  Thus every point where \(\GG\) has non-zero
stalk is a point where \(\FF\) has non-zero stalk.
\end{solution}

\begin{exercise}[Invertible sheaves are locally free rank one]\label{ex:line-bundle-local}
Show that an \(\OO_X\)-module \(\LL\) is invertible if and only if it is locally
free of rank one.
\end{exercise}

\begin{solution}
This is a matter of terminology in one direction: invertible was defined by
local isomorphisms \(\LL|_U\simeq\OO_U\), which is local freeness of rank one.
Conversely, a locally free rank-one sheaf has dual locally isomorphic to
\(\OO_U\), and the evaluation map \(\LL\tensor\LL^\vee\to\OO_X\) is locally the
identity, hence globally an isomorphism.  Thus it is invertible in the Picard
group.
\end{solution}

\begin{exercise}[Closed subschemes from ideals]\label{ex:ideal-sheaf-affine}
Let \(X=\Spec A\) and \(I,J\subset A\).  Show that
\[
    \Spec(A/I)\subset \Spec(A/J)
\]
as closed subschemes of \(X\) if and only if \(J\subset I\).
\end{exercise}

\begin{solution}
The inclusion of closed subschemes corresponds to a factorization
\[
    A\to A/J\to A/I.
\]
Such a factorization exists if and only if every element of \(J\) maps to zero
in \(A/I\), i.e. \(J\subset I\).  This reverses inclusions because larger ideals
define smaller closed subschemes.
\end{solution}

\section{Serre Correspondence on Projective Schemes}

\begin{definition}
Let \(S=A[x_0,\ldots,x_n]\) and \(X=\Proj S\).  A graded \(S\)-module \(M\) is
\emph{irrelevant torsion} if for every \(m\in M\), some power of the irrelevant
ideal \(S_+=(x_0,\ldots,x_n)\) annihilates \(m\).
\end{definition}

\begin{proposition}
If \(M\) is irrelevant torsion, then \(\widetilde M=0\) on \(\Proj S\).
\end{proposition}

\begin{proof}
It is enough to check on the standard opens \(D_+(x_i)\).  If \(m\in M\), then
\((S_+)^N m=0\) for some \(N\), so in particular \(x_i^Nm=0\).  Hence \(m/1=0\)
in \(M_{x_i}\), and therefore the degree-zero localization \(M_{(x_i)}\) is
zero.  Thus all sections on the standard affine cover vanish.
\end{proof}

\begin{theorem}[Serre correspondence, projective space form]
Let \(A\) be noetherian and \(X=\PP^n_A\).  Every coherent sheaf on \(X\) is of
the form \(\widetilde M\) for some finitely generated graded
\(A[x_0,\ldots,x_n]\)-module \(M\).  Two such modules define isomorphic sheaves
if and only if they become isomorphic after quotienting by irrelevant torsion
and ignoring finitely many low degrees.
\end{theorem}

\begin{proof}
Let \(\FF\) be coherent.  By Serre global generation, for \(m\gg0\), the sheaf
\(\FF(m)\) is generated by finitely many global sections.  These sections give
a surjection
\[
    \OO_X(-m)^{\oplus r}\to\FF.
\]
The kernel is coherent; applying the same argument to the kernel and iterating
once gives a finite presentation by sums of twists of \(\OO_X\).  Translating
this presentation into graded free modules over \(S=A[x_0,\ldots,x_n]\) gives a
finite graded module \(M\) with \(\widetilde M\simeq\FF\).  The ambiguity is
exactly irrelevant torsion and finite changes in low degree, because Proj only
sees degree-zero localizations at homogeneous elements of positive degree.
\end{proof}

\section{Fitting Ideals}

\begin{definition}
Let \(M\) be a finitely presented module over a ring \(A\), with presentation
\[
    A^m\xto{\Phi} A^n\to M\to0.
\]
The \(r\)-th \emph{Fitting ideal} \(\operatorname{Fitt}_r(M)\) is generated by
the \((n-r)\times(n-r)\) minors of the matrix \(\Phi\), with the conventions
\(\operatorname{Fitt}_r(M)=A\) for \(r\ge n\) and \(0\) if \(n-r>m\).
\end{definition}

\begin{fact}
Fitting ideals are independent of the chosen finite presentation and commute
with localization and base change.
\end{fact}

\begin{definition}
If \(\FF\) is a coherent sheaf on a locally noetherian scheme \(X\), the
\emph{Fitting ideal sheaf} \(\operatorname{Fitt}_r(\FF)\subset\OO_X\) is
obtained by applying the module construction on affine opens.
\end{definition}

\begin{proposition}
The locus
\[
    \{x\in X:\dim_{\kappa(x)}\FF\tensor\kappa(x)\le r\}
\]
is open, and is the complement of the closed subset defined by
\(\operatorname{Fitt}_r(\FF)\).
\end{proposition}

\begin{proof}
The assertion is affine local.  Let \(X=\Spec A\) and let \(M\) be presented by
a matrix \(\Phi:A^m\to A^n\).  After tensoring with \(\kappa(x)\), the fiber has
dimension at most \(r\) exactly when the rank of \(\Phi(x)\) is at least
\(n-r\).  This is equivalent to at least one \((n-r)\times(n-r)\) minor being
non-zero at \(x\), i.e. to \(x\notin V(\operatorname{Fitt}_r(M))\).  Hence the
locus is open.
\end{proof}

\section{Flatness and Fibers of Coherent Sheaves}

\begin{fact}[Local criterion for flatness]
Let \((A,\mathfrak m)\to(B,\mathfrak n)\) be a local homomorphism of
noetherian local rings, and let \(M\) be a finite \(B\)-module.  Then \(M\) is
flat over \(A\) if and only if
\[
    \Tor_1^A(A/\mathfrak m,M)=0
\]
and the expected lifting of relations modulo \(\mathfrak m\) holds.  In
particular, if \(A\) is a field, every \(M\) is flat over \(A\).
\end{fact}

\begin{proposition}
Let \(f:X\to S\) be locally of finite type with \(S\) locally noetherian, and
let \(\FF\) be coherent on \(X\).  The set of points \(x\in X\) where \(\FF_x\)
is flat over \(\OO_{S,f(x)}\) is open when \(\FF\) is flat at all generic points
of the fibers.
\end{proposition}

\begin{proof}
The question is local on \(X\) and \(S\), so reduce to a finite module \(M\) over
a finite type algebra \(B\) over a noetherian ring \(A\).  Generic flatness says
that after inverting a suitable element of \(A\), the module \(M\) becomes flat
over \(A\).  Applying this statement on the closures of points of \(S\) and
using noetherian induction gives an open flatness locus.  The hypothesis on
generic points ensures that the open neighbourhoods obtained in this way cover
the desired locus.
\end{proof}
